\documentclass[12pt]{article}
\usepackage[T1]{fontenc}
\usepackage[utf8]{inputenc}
\usepackage{lmodern}
\usepackage{textcomp}
\usepackage{enumitem}
\usepackage{geometry}
\geometry{margin=1in}
\begin{document}
\begin{center}
Answer Key - Mathematics - Graduate Level Assessment 8 \
\small (Answers are ordered exactly as in the question paper.)
\end{center}

\section*{Multiple Choice}
\begin{enumerate}[label=\arabic*.]
\item \textbf{Answer:} A
\\ \textit{Options:} A) 0.5; B) 1; C) 0.95; D) 0.65; E) 0.85
\\ \textit{Note:} The correct answer is 0.5 because, given that there is exactly one arrival in the interval [0, t], the first arrival time \(X_1\) is uniformly distributed over that interval, leading to \(P(X_1 \leq t/2) = \frac{t/2}{t} = 0.5\).
\item \textbf{Answer:} E
\\ \textit{Options:} A) 0.8; B) 0.3; C) 0.1; D) 0.5; E) 0.4
\\ \textit{Note:} The correct answer is 0.4 because the solution to the initial value problem, which is a second-order linear homogeneous differential equation, can be determined by finding its characteristic polynomial, solving for the roots, and using the initial conditions to construct the explicit solution, which reveals that the first instance where the solution equals zero occurs at t = 0.4.
\item \textbf{Answer:} A
\\ \textit{Options:} A) False, False; B) True, True; C) Not Determinable, True; D) Not Determinable, False; E) True, False
\\ \textit{Note:} The correct answer is false for both statements because a factor group of a non-Abelian group can be Abelian, and while normality is transitive under certain conditions, K being a normal subgroup of H and H being a normal subgroup of G does not guarantee that K is normal in G unless H is also a normal subgroup of K.
\item \textbf{Answer:} C
\\ \textit{Options:} A) \$4.52; B) \$7.28; C) \$5.52; D) \$6.28; E) \$3.72
\\ \textit{Note:} The correct answer is $5.52 because the cost per stamp is $8.28 divided by 18, which equals approximately $0.46 per stamp, and multiplying this rate by 12 stamps gives a total of $5.52.
\item \textbf{Answer:} C
\\ \textit{Options:} A) -1; B) -13; C) -4; D) -3; E) 0
\\ \textit{Note:} The slopes of two perpendicular lines are negative reciprocals of each other, and since the slope of the line represented by $x + 2 y + 3 = 0$ is -1/2, the slope of the line represented by $ax + 2 y + 3 = 0$ must be 2, leading to the equation $-\frac{a}{2} = 2$, which simplifies to $a = -4$.
\end{enumerate}

\section*{True / False}
\begin{enumerate}[label=\arabic*.]
\setcounter{enumi}{5}
\item \textbf{Answer:} True
\\ \textit{Options:} A) True; B) False
\\ \textit{Note:} The statement is true because the given conditions imply that the sequence \( a_n \) is bounded between 0 and 1 and converges to a limit, and the inequality \( (1-a_n)a_{n+1} > 1/4 \) ensures that \( a_{n+1} \) approaches a value close to 0.5 as \( n \) increases, thereby establishing that the limit of the sequence as \( n \) approaches infinity is indeed 0.5.
\item \textbf{Answer:} False
\\ \textit{Options:} A) True; B) False
\\ \textit{Note:} The expression -25 - (-11) simplifies to -25 + 11, which equals -14, thus making the statement false.
\item \textbf{Answer:} True
\\ \textit{Options:} A) True; B) False
\\ \textit{Note:} The statement is true because in a nonzero free abelian group, any independent set of generators can be extended to form a new basis, leading to the existence of infinitely many distinct bases.
\item \textbf{Answer:} True
\\ \textit{Options:} A) True; B) False
\\ \textit{Note:} The determinant of matrix A can be calculated using the formula for the determinant of a 3 x 3 matrix, which yields a value of 200 when the elements of the matrix are substituted into the formula.
\item \textbf{Answer:} True
\\ \textit{Options:} A) True; B) False
\\ \textit{Note:} The difference of 142.76 and 16.5 is accurately calculated as 126.26 because when subtracting 16.5 from 142.76, the result is indeed 126.26, confirming the correctness of the arithmetic.
\end{enumerate}

\section*{Long Form Response}
\begin{enumerate}[label=\arabic*.]
\setcounter{enumi}{10}
\item \textbf{Answer:} The length of the hypotenuse is given by the distance formula to be $\sqrt{(5-(-1))^2 + (-5-(-1))^2} = \sqrt{6^2+4^2} = \sqrt{52}$. The length of the leg is then given by $\sqrt{52}/\sqrt{2} = \sqrt{26}$ (alternatively, the Pythagorean Theorem can be applied), and the area of the isosceles right triangle is then equal to $\frac 12 \cdot \sqrt{26} \cdot \sqrt{26} = \boxed{13}.$
\\ \textit{Note:} correct because it accurately calculates the length of the hypotenuse using the distance formula, derives the length of the legs of the isosceles right triangle from the hypotenuse, and correctly applies the formula for the area of a triangle to find the area of triangle ABC.
\item \textbf{Answer:} To determine the maximum possible value of the quotient \(\frac{x}{y}\) for distinct elements \(x\) and \(y\) from the set \(\{\frac{2}{5}, \frac{1}{2}, 5, 10\}\), we analyze the possible pairings. The largest value of \(x\) is \(10\) and the smallest value of \(y\) is \(\frac{2}{5}\). Thus, the maximum quotient is \(\frac{10}{\frac{2}{5}} = 10 \times \frac{5}{2} = 25\). Checking other combinations, such as \(\frac{5}{\frac{1}{2}} = 10\) and \(\frac{5}{\frac{2}{5}} = 12.5\), confirms that none exceed \(25\). Therefore, the maximum possible value of \(\frac{x}{y}\) is \(25\).
\\ \textit{Note:} correct because the maximum quotient \(\frac{x}{y}\) is achieved by selecting the largest element \(x = 10\) and the smallest element \(y = \frac{2}{5}\), resulting in a value of 25, which is greater than any other possible pair of distinct elements from the set.
\end{enumerate}

\vfill
\noindent\textit{End of Answer Key}
\end{document}